\section{Discussion and open problems}

The former assertion that FFF squares of order $n>1$ exist exactly
at powers of two is false: the displayed order-$12$ witness is a
counterexample. Nontrivial odd orders, order $6$, and order $10$
are excluded; the order-$10$ argument is an explicit computational
dependency imported from the full companion report. Order $18$
remains the next undecided small even non-power-of-two order.

The row-fibred theorem preserves the exact three-view pattern even
with different fibres. Under the additional binary-quotient-free
hypothesis, Theorem~\ref{thm:affine-fibre-orbits} classifies the
fixed elementary-abelian-quotient family by affine orbits.
Computational Theorem~\ref{cthm:fff20-1703} gives lower
bounds for all FFF classes, not exact censuses or unrestricted
direct-product cancellation.

The next mathematical questions are:
\begin{enumerate}
\item Decide unrestricted FFF existence at order $18$ and sharpen
the minimal dyadic-exponent function $h(m)$.
\item Find structural obstructions coupling all three views that
do not also exclude the explicit order-$12$, $14$, and $98$ controls.
\item Determine what lies outside the fixed constructive families
and whether quotient recovery survives weaker hypotheses.
\end{enumerate}

The full report records several exact, but restricted, order-$18$
construction exclusions. None covers all order-$18$ Latin squares.
Timeouts in the smaller equivalent coloring model are likewise
not exclusions. No global Alon--Tarsi conclusion is claimed
\cite{AlonTarsi1992}; literature priority and independent expert
review remain open.
